% Layer 3 — Inference Architectures (draft pass, 2026-10-17).
% Insertion point: \section{Layer 3 --- Inference Architectures}\label{sec:architectures} in main.tex.
% Cite policy: bib keys where they exist; otherwise placeholder keys of the form p2[<verified arXiv ID>].

\subsection{Generalists, Specialists, and Hybrid Pipelines}
\label{sec:l3-mix}

Layer~3 covers the choices that govern how a trained model is \emph{used}: at what granularity, with what conditional computation, and with which external calls. The first choice is granularity. A generalist is asked to produce everything at once --- answer, format, and whatever confidence the response needs; a specialist commits to narrower outputs that a downstream system can verify mechanically: boxes, discrete labels, bounded scores. That verifiability is the argument for the specialist architecture, and also the one most often asserted rather than measured end to end. \evvendor{Interfaze positions an orchestrated set of task-specific small models against general-purpose incumbents \cite{interfaze2026}}; the claim lives at component level, and \cref{fig:case-hybrid} reconstructs one such pipeline from vendor descriptions chiefly so that the edges the vendor leaves undocumented --- monitors, validators, escalation --- stay visible in our account. The hybrid pays for composition: every added component brings its own error distribution, and the pipeline's end-to-end reliability depends on how those errors propagate, not on the best component's headline number. We therefore treat hybrid pipelines as stack instances with their own failure envelopes rather than as upgraded single models.

\subsection{Mixture-of-Experts: Routing Inside the Model}
\label{sec:l3-moe}

Mixture-of-Experts is conditional computation inside the model: capacity scales with parameter count while cost scales with activated experts, and the router becomes a learned, partially inspectable dispatch policy \cite{id240706204}. Our selection framework needs the intra-model and system-level split more sharply than the survey literature draws it, so we keep two routing loci strictly apart: \emph{intra-model} routing decides which experts fire for a token and is an architecture property; \emph{system-level} routing decides which model or cascade tier receives a request (\cref{sec:l3-routing}) and is an orchestration policy. On reliability, MoE-RBench is the benchmark evidence we anchor: \evindep{Across safety and hallucination, adversarial resilience, and out-of-distribution robustness, sparse MoE can be made more reliable than dense models of its class, but robustness is sensitive to basic training settings and hyperparameters \cite{id240611353}}. Reliability here is a tunable, configuration-dependent property. Compression follows the same logic: \evindep{structured expert pruning preserves in-domain utility at moderate ratios, whereas hallucination risk rises at extreme pruning and both utility and factual reliability degrade quickly out of domain \cite{id260701444}}; at deployment time compression is a reliability decision as much as a memory one.

\subsection{Routing and Cascades: Cost--Quality Allocation as an Architectural Lever}
\label{sec:l3-routing}

Above the model, allocation happens per query. FrugalGPT established the cascade pattern (chains of increasingly capable models whose cheaper members are trusted only when their scored responses clear a cutoff), reallocating spend toward the queries that need it \cite{frugalgpt2023}. The architectural weak point is the dispatch signal: cascades greedily select the lightest model from confidence scores, but deep models are systematically overconfident and confidence distributions differ across models and languages, so the C3 system \cite{id240215991} interposes a per-member calibration stage before escalation. SATER unifies pre-generation routing with cascade-style escalation under a token-efficiency objective, targeting the accuracy--latency split between the two families \cite{id251005164}. Routing also extends beyond a fixed pool: UniRoute routes to previously unseen models by representing each candidate through its predictions, with excess-risk bounds relating the learned rule to a theoretical optimum \cite{id250208773}. The evidence base remains fragmented. The characterization survey of routing and hierarchical inference is efficiency-scoped and notes the absence of a unified evaluation across methods \cite{id250606579}, and the risk-keyed selection question stays open (\cref{sec:selection}).

\subsection{Looped and Recurrent Depth: Compute Scaling Without Parameter Scaling}
\label{sec:l3-depth}

Depth is the second axis, and it can be repeated rather than stored: looped transformers apply one compact stack for multiple rounds, buying unrolled depth at fixed parameter count. \evtheory{DeepLoop's tied-depth analysis yields a first-order perturbation bound controlled by a visit-alignment coefficient --- the shared weight update aggregates gradients from repeated visits --- recovering standard residual-scaling exponents when visits decorrelate and requiring larger ones in the aligned regime} \cite{id260713491}. \evaurep{The DeepLoop authors' accompanying ablations are small-scale (GPT-style looped language models), effectively tied with the unlooped baseline at one loop and improving at moderate loop counts, but single-seed}, so we record a stability envelope rather than a performance verdict. Engineering follow-ups attack the family's latency cost with self-speculative decoding \cite{id260917184,id261021594}. The broader recurrent lineage makes the same wager with fixed-size state: the RWKV survey documents the merge of recurrence and attention \cite{id241214847}, and a controlled 72-model study of hybrid linear attention finds that \evindep{standalone superiority of a linear-attention variant does not transfer to its hybrids, with recall governed by the linear-to-full mixing ratio} \cite{id250706457}, so the envelope belongs to the hybrid rather than to either component.

\subsection{Memory, Adaptation, and Latent Reasoning at Inference Time}
\label{sec:l3-memory}

The same reuse principle also operates inside the weights, where it yields memory learned at test time: Titans adds a neural long-term memory that memorizes historical context while attention handles the current window \cite{id250100663}, and TPTT retrofits pretrained transformers with linearized attention and memory gating without full retraining \cite{id250617671}. Adaptation at inference then stops being a deployment-time constant: \evaurep{verifier-driven test-time training selects pseudo-labeled examples through a learned verifier and adapts the model out of distribution} \cite{id250519475}, which is why this row lives only partly on Layer~3; the verifier's own reliability, the drift adaptation induces, and the monitoring it requires are assurance concerns (\cref{sec:assurance}). Hierarchical recursive reasoning pushes the same logic to continuous depth: HRM demonstrated strong results on 2D reasoning tasks with compact latent recurrence rather than pretraining scale \cite{id250621734}, and the evidence-discipline lesson lies in the delta between those demos and the general-reasoning claims they invite. A critical review re-examines the design space, tests variants, and corrects misconceptions of an explicitly early-stage family \cite{id251000355}, while a dynamical-systems reformulation recasts the recursion as continuous dynamics with an explicit equilibrium requirement and feature-collapse mitigation, addressing the instability the original training dynamics do not guarantee away \cite{id260322871}. Latent reasoning also narrows the controls we take for granted in token space: \evindep{activation steering of continuous thoughts produces substantially weaker effects on generation than steering explicit chain-of-thought, even at comparable movement inside the representation}, a latent-to-language transition gap \cite{id260921662}, and detection work builds benchmarks to probe misaligned reasoning planted inside continuous thought \cite{id260423460}.

\subsection{Tool-Augmented Agents}
\label{sec:l3-agents}

The final move re-opens the model to the world: ReAct interleaves reasoning traces with actions that query external environments, letting observations steer subsequent reasoning \cite{id221003629}. Architecturally, each tool call is a verification opportunity --- deterministic functions, retrieval, and execution are checkable in ways token generation is not (\cref{sec:contracts,sec:assurance}) --- and simultaneously a new failure path: planner errors propagate through execution, intermediate state persists, and some actions are irreversible. The agent surveys chart precisely this surface, factoring planning into task decomposition, plan selection, external modules, reflection, and memory \cite{id240202716}, and organizing broader agent methodology around architecture, collaboration, and evolution \cite{id250321460}. None of these taxonomies carries a reliability guarantee; they enumerate where validation must be placed, and because they relocate rather than remove the reliability burden, they make the strongest case in this layer for co-designing L3 and L5.

\subsection{What Layer 3 Does Not Guarantee}
\label{sec:l3-envelope}

Read as a whole, the layer offers a set of allocation choices: where to spend parameters (MoE), repeated depth (looping), or calls (cascades). Each of those mechanisms is a candidate whose envelope we have tried to state. None of them makes the model calibrated by construction. Cascades and routers consume confidence estimates that base models oversupply and miscalibrate \cite{id240215991}, compressed experts shift where hallucination risk sits \cite{id260701444}, and test-time memory moves the training--inference boundary into deployment \cite{id250519475}. End-to-end reliability remains a system property, and component-level wins --- even independent component benchmarks such as the decision-tier evaluations of the Jev System-One family \cite{rauk2026jevbenchmark} --- do not compose into system guarantees. Layer~3 supplies candidates and their costs; what the system may claim is decided by the selection procedure and the assurance layer (\cref{sec:selection}, \cref{sec:assurance}).
